\scanpage{25}\mspage{24}
Finalement, on travaillera avec la top. définie d'abord (base formée avec les ``intervalles''), et on appellera néanmoins par abus de langage, les $[a,b]$, $\unsure{L|\mathscr{L}}_{\leq a}$, $L_{\geq b}$ (ainsi que $L$) \uline{sous-tronçons} ordinaires (ou \unsure{biordonnés|bien donnés|bien disons|binominés}) ``fermés'', [NB on les suppose \unsure{non|ici} réduits à un pt, i.e.
\[
\deleted{\unsure{[|i}} a<b\quad \text{resp.}\quad a\ne\mathrm{or}(L)\quad \text{resp.}\quad \deleted{\unsure{a}}\unsure{b|a}\ne\mathrm{ex}(L)\;]
\]

Un \uline{sous-tronçon \inserted{fermé} large} est une \deleted{sous-tronçon} partie qui est soit un \unsure{ss-tronçon|s-tronçon}, soit \unsure{réduit|réduite} à un pt.

Pour une telle partie, $T$, on
% Widely spaced strokes cross the whole lower block, from pose to the last example.
% The close zigzags and local horizontal cancellations show an earlier level of correction.
\deleted{pose}
\[
\begin{aligned}
\deleted{\partial T=}&\deleted{\{t\in T\mid t\in \overline{L\setminus T}\}}\quad\text{\deleted{bord}}\\
 &\text{\deleted{\unsure{frontière dans la top.|frontière dans la top}}}\\
 &\text{\deleted{absolu de $T$ $=$ ensemble}}\\
 &\text{\deleted{\unsure{canonique}}}\\
 \deleted{={}\unsure{T\cap L|T\cap\overline{L}}}&\quad\text{\deleted{fermé de ses extrémités}}\\
 &\deleted{(\mathrm{NB}\ \mathrm{card}\,\partial T\in\{0,1,2\})}
\end{aligned}
\]

\deleted{On voit que pour un tronçon fermé (\uline{NB} non réduit à un pt) on a}

\[
\begin{aligned}
 \deleted{\partial[a,b]=\{a,b\}}&\qquad \deleted{\partial L_{\leq a}=\{a\},\quad \partial L_{\geq a}=\{a\}}\\
 &\deleted{\partial L=\emptyset}
\end{aligned}
\]
